\documentclass[10pt]{article}
\usepackage{icmlko}

\icmlmeta{IP 파운데이션 월드모델}{971}{연관에서 예측으로}{아홉 사이클의 모든 $\rho$ 는 같은 자료 안의 연관이었다 --- 유보에서 재니 들뜸은 예측을 $+0.0255$ 돕고, 그 값은 채택 조건 다섯 중 넷만 넘는다}

\begin{document}

\twocolumn[
\icmlheader

\begin{abstract}
이 연구 줄기는 아홉 사이클(962--970) 동안 \emph{「집단 관심의 기억은 90일보다 길다」}는 명제를
편 스피어만 $\rho$ 와 순열 $p$ 로 재 왔다. \textbf{그 모든 수는 같은 자료 안의 연관이다} ---
학습 행과 유보 행을 가르지 않았다. 두 차례의 외부 비평이 연달아 같은 말을 했다:
\emph{「예측 성능은 한 번도 안 쟀다.」}
이 노트는 자를 뗀다. 같은 규격의 자료를 판이 쓰는 그 시간 분할($T = 2025.0$)로 갈라
\textbf{학습 행에서만 적합하고 유보 행의 결과를 예측}시킨 뒤,
\textbf{그 예측과 실제의 스피어만}을 잰다. 통제만 쓴 팔 $B$ 와 통제에 들뜸을 더한 팔 $C$ 의
차가 자다. 🔴 \textbf{측정 전에 채택 조건 다섯을 못 박았고}, 이 자로 잴 수 있는 최소 효과와
\textbf{이 자로 직접 잰 잡음 바닥}(무작위 열 50 뽑기)을 같이 냈다.
결과: 도메인 \textbf{7} · 유보 행 \textbf{532} 에서
$\Delta\rho_{\mathrm{pred}} = \mathbf{+0.025527}$ 이고 \textbf{7 도메인 전부에서 양수}이며
저장소 잡음 바닥 $0.01055$ 와 이 자로 잰 바닥 $0.016593$ 을 \textbf{둘 다 넘는다}.
\textbf{그러나 짝 붓스트랩 95\% 구간이 $[-0.0163,\ 0.0702]$ 로 0 을 품는다} ---
사전등록 채택 조건 \textbf{다섯 중 넷}만 넘었다. 🔴 \textbf{이것은 「이 자를 못 넘었다」이지
「명제가 거짓이다」가 아니다.} 곁으로, 같은 유보 행에서 잰 \emph{연관} $\rho$(유보 행 가중 $|\rho| = 0.145443$)는
예측 $\Delta$ 의 \textbf{5.7 배}다 --- \textbf{「연관이 있다」에서 「예측을 돕는다」로 바로 못 건너간다는 것이
처음으로 수가 됐다.}
\end{abstract}
]

\section{자를 왜 뗐나}

앞선 아홉 사이클은 도메인마다 개체 집합을 하나 잡고,
\emph{들뜸}($x$ = 개시일 직전 90일 로그 일평균 $-$ 그 앞 365일 로그 일평균)과 결과 $y$ 의
편 스피어만을 통제 둘($\texttt{wiki\_level}$ · $\texttt{recent\_term}$) 아래에서 재고
순열로 $p$ 를 냈다. \textbf{그 절차는 학습 행과 유보 행을 안 가른다.} 그래서 나오는 수는
정의상 \textbf{같은 자료 안의 연관}이고, \emph{「월드모델」}이라는 말이 요구하는
\emph{「못 본 것을 맞힌다」}에는 답하지 않는다.

\claimx{자를 유보 예측으로 바꾸는 것은 통제 선택의 순환을 끊는 유일한 길이다}{학습 행에만 적합했는데도 유보 채점이 학습 통계에 의존하면 --- 즉 유보 결과를 통째로 갈았을 때 예측이 바뀌면 --- 이 주장은 떨어진다.}

앞 노트의 비평이 명제의 바닥에 관해 이렇게 적었다:
\emph{「되살린 통제로 판이 $-0.002962$ 내려갔다는 건 「그 상수가 예측을 돕는다」는 뜻이라,
명제는 앞으로도 「유보 정답으로 채점한 판 $\rho$ 를 보고 고른 통제 집합」 위에 선다.」}
\textbf{이 순환은 자를 바꾸지 않으면 안 끊긴다.} 유보에서 재면 통제 열은
\textbf{학습 행에서만} 가짓수 게이트를 통과하고 \textbf{학습 행에서만} 계수를 얻으며,
채점은 그 적합이 \textbf{한 번도 못 본 행}에서 난다.

배선 검사가 그것을 실측으로 확인한다. 적합기 \texttt{fit\_predict(Xtr, ytr, Xho)} 의
\textbf{인자에 $y_{\mathrm{ho}}$ 가 없고}, 유보 결과를 통째로 갈아 넣어도 예측의
최대 절대차가 \textbf{$0.0$} 다(허용 $10^{-12}$). 배선 아홉 자리 전부 초록이다(\textbf{9 / 9}).

\section{검출력 줄 --- 측정 전에 채웠다}

\begin{center}\tsize
\begin{tabular}{lrrr}
\toprule
도메인 & 학습 & \textbf{유보} & 최소 효과(2 SE) \\
\midrule
세계애니 & 845 & \textbf{230} & 0.1322 \\
애니 & 248 & \textbf{114} & 0.1881 \\
게임 & 93 & \textbf{52} & 0.2801 \\
시장팝업 & 41 & \textbf{42} & 0.3123 \\
웹툰 & 158 & \textbf{38} & 0.3288 \\
모바일 & 120 & \textbf{29} & 0.3780 \\
도서 & 20 & \textbf{27} & 0.3922 \\
\midrule
\textbf{합} & \textbf{1{,}525} & \textbf{532} & 묶음 \textbf{0.0873} \\
\bottomrule
\end{tabular}
\end{center}

판 12 도메인 중 \textbf{7} 이 「학습 $\ge 20$ · 유보 $\ge 20$」을 채운다.
빠진 다섯의 사유를 그대로 적는다: 팝업 학습 \textbf{18}(문턱에 \textbf{둘} 모자란다) ·
아이돌 유보 6 · 만화 유보 \textbf{0} · 영화·펀딩 규격 행 \textbf{0}·\textbf{2}.
🔴 \textbf{판 유보 전량 3{,}775 중 이 규격이 닿는 유보는 573(15.2\%)이다.}

\section{결과}

\claimx{들뜸을 통제에 더하면 유보 예측이 $+0.025527$ 좋아진다 --- 7 도메인 전부에서 양수이고, 두 잡음 바닥을 둘 다 넘는다}{무작위 열 50 뽑기의 95 분위(이 자로 직접 잰 잡음 바닥)가 관측 $\Delta$ 이상이면 떨어진다. 실측 0.016593 $<$ 0.025527.}

\begin{center}\tsize
\begin{tabular}{lrrrr}
\toprule
도메인 & 유보 & $\rho_B$ & $\rho_C$ & $\Delta$ \\
\midrule
시장팝업 & 42 & 0.0669 & 0.2054 & \textbf{$+$0.1385} \\
애니 & 114 & 0.0290 & 0.0682 & $+$0.0392 \\
웹툰 & 38 & $-$0.2028 & $-$0.1726 & $+$0.0302 \\
도서 & 27 & $-$0.0180 & 0.0107 & $+$0.0287 \\
게임 & 52 & 0.4947 & 0.5088 & $+$0.0141 \\
세계애니 & 230 & 0.4358 & 0.4383 & $+$0.0025 \\
모바일 & 29 & 0.7011 & 0.7028 & $+$0.0017 \\
\midrule
\textbf{묶음} & \textbf{532} & --- & --- & \textbf{$+$0.025527} \\
\bottomrule
\end{tabular}
\end{center}

\claimx{그러나 채택 조건 다섯 중 넷만 넘었다 --- 짝 붓스트랩 구간이 0 을 품는다}{구간의 아래끝이 0 을 넘으면 이 주장은 떨어진다. 실측 $[-0.016265,\ 0.070222]$ · 짝 SE 0.021669 · $|\Delta| \div$ SE $= 1.18$.}

측정 \textbf{전에} 못 박은 다섯은 이렇다. ① $\Delta > 0$ \textbf{✓} ·
② $\Delta \ge 0.01055$(저장소 잡음 바닥) \textbf{✓} ·
③ $\Delta \ge$ 이 자로 잰 바닥 $0.016593$ \textbf{✓} ·
④ 짝 붓스트랩 95\% 구간 아래끝 $> 0$ \textbf{✗} ·
⑤ 7 중 $\ge 5$ 도메인 양수 \textbf{✓}(실측 \textbf{7 / 7}).
🔴 \textbf{④에서 걸렸다.} 그러므로 결론은 \emph{「이 자를 못 넘었다」}이고,
\textbf{「명제가 거짓이다」가 아니다} --- 이 표본에서 짝을 안 지었을 때 가를 수 있는
최소 효과가 묶음 \textbf{0.0873} 이고, 그보다 작은 참 효과는 이 자가 원리상 못 가른다.

🔴 \textbf{정직 신고}: 붓스트랩은 \textbf{유보 행만} 되뽑는다. \textbf{학습 재표집은 안 했다.}
그러므로 이 폭은 \emph{「적합을 고정했을 때」}의 폭이고 적합의 변동은 안 들어 있다.

\section{연관과 예측은 종류가 다르다}

\claimx{같은 유보 행에서 잰 연관 $\rho$ 는 예측 $\Delta$ 보다 여러 배 크다}{유보 행 가중 $|\rho_{\text{연관}}|$ 이 $\Delta\rho_{\mathrm{pred}}$ 이하이면 떨어진다.}

같은 행에서 \textbf{연관}(들뜸 $\leftrightarrow$ 결과)을 재면 도메인별로
$|\rho|$ 가 $0.076$--$0.320$ 이고 유보 행 가중 평균이 \textbf{0.145443} 이다. \textbf{예측 $\Delta$ 는 $0.0017$--$0.1385$ 이고 묶음이 $0.0255$ 다 --- 연관이 \textbf{5.7 배} 크다.}
그리고 \textbf{들뜸만 쓴 팔 $A$ 는 통제만 쓴 팔 $B$ 보다 묶음에서 $-0.1256$ 나쁘다} ---
\textbf{들뜸은 통제를 대체하지 못하고 보탤 뿐이다.}
🔴 세 도메인(도서 · 세계애니 · 시장팝업)에서는 \textbf{학습에서 배운 부호가 유보에서 뒤집힌다}
--- 팔 $A$ 의 $\rho$ 부호가 유보 연관의 부호와 반대다. \textbf{연관은 있는데 방향이 안 건너간다.}

\section{자를 자기 러너에 걸었다}

\claimx{항진명제 계수는 리팩터링만으로 내려간다 --- 심어서 보였다}{인라인으로 심은 항진명제와 같은 값을 인자로 옮긴 항진명제를 자가 똑같이 잡으면 떨어진다.}

앞 노트는 \texttt{out970\_meta.json} 을 \textbf{안 냈다}. 이 노트는 낸다 ---
그리고 분모에 \textbf{자기 러너와 자 자신을 둘 다} 넣었다.
비평이 대신 걸어 낸 밑값(\texttt{recommit970}: 분모 6 · 항진명제 0 ·
\textbf{모른다 5 (83.3\%)} · 증명된 낙하 1 (16.7\%))을 \textbf{비트로 재현}했고,
이 노트의 러너는 \textbf{분모 13 · 항진명제 0 · 모른다 1 · 증명된 낙하 4 (30.8\%)} 다.

🔴 \textbf{그러나 자랑하면 안 된다.} 같은 항진명제를 두 꼴로 심어 재 봤다(분모 4):
\textbf{인라인} 꼴은 항진명제로 \textbf{잡히고}, \textbf{비트로 같은 값을 인자로 넘긴} 꼴은
\textbf{「상수 False = 모른다」로 샌다.} 게다가 자는 \textbf{그 죽은 검사와 진짜 검사를 못 가른다}
--- \textbf{둘 다 「상수 False」다.} 그러므로 이 노트가 앞 노트의 세 사이클 묵은 항진명제
(\texttt{W6})를 \textbf{1 → 0} 으로 만든 것은 \emph{자의 계수}로는 증거가 아니다.
증거는 다른 데 있다: 고친 \texttt{W6} 은 \textbf{양성 대조}(같은 바이트의 sha 대
한 줄 더한 바이트의 sha)를 먹으므로, \textbf{sha 를 잘라 비교하는 자로 바꾸면 그 자리에서
떨어진다.} \textbf{구판에는 그 갈래가 원리상 없었다.}

\section{앞 노트의 수 정정}

\claimx{앞 노트에서 「없어진 것」은 둘이 아니라 셋이고, 그 틀린 수가 상시 조항 정본에 박혀 있었다}{diff 가 내용을 내는 자리가 셋이 아니면 떨어진다.}

셋이다: \textbf{(i)} \texttt{audit\_cells()} 의 \emph{「조항 60」} 키 ·
\textbf{(ii)} \texttt{history()} 의 \S H 사료 · \textbf{(iii)} \texttt{compare()} 의 \S E.
🔴 \textbf{앞 노트 자신의 사전등록이 \S E 를 찾기도 전에 (i)(ii) 를 지목했고 자기 원장도 셋을
세어 적었는데}, 「둘」이 \textbf{다섯 자리}에 박혔고 그중 하나가 \textbf{상시 조항 정본}이었다.
\textbf{앞선 비평이 「하나」라 해서 「둘」로 고친 것이라} ---
\textbf{남이 준 수를 그대로 받아 고치면 같은 실패가 한 칸 옮겨서 반복된다.}
넷을 고쳤고 다섯째(커밋 메시지)는 \textbf{못 고친다}.

곁으로: 잎 분모 \textbf{3{,}644 는 산출물 「다섯」의 것}이고 여섯째에는
\textbf{줄번호도 sha 도 시각도 아닌 진짜 차이}(연 수 $2{,}241 \to 2{,}154$)가 있다 ·
파수병 분모는 \textbf{14 / 33 잴 수 있는 것 $+$ 1 못 재는 것}이다
(\texttt{nan\_to\_num} 은 파수병이 원리상 못 센다) ·
카드의 \emph{「합친 sha」}는 \textbf{어떤 인자에서도 안 나오는 수}라 도장이 아니어서 지웠다.

\section{못 한 것}

채택 조건 ④를 못 넘었다(유보 532 행이 이 자의 분해능이다) ·
붓스트랩이 학습을 재표집 안 한다 ·
\emph{「상수 False」} 안에서 죽은 검사와 진짜 검사를 가르는 자를 \textbf{못 만들었다}
--- \textbf{못 가른다는 것만 보였다} ·
헤드라인 도메인(팝업)이 학습 \textbf{18} 행이라 예측 팔에 못 들어갔다.
🔴 \textbf{측정 뒤에 문턱을 옮기지 않았다} --- 그것이 사후 부분집합이다.

\end{document}
